\chapter{Khi Nào Em Mới Thôi Sợ Đi Học}
\markboth{Khi Nào Em Mới Thôi Sợ Đi Học}{When Will I Stop Being Afraid of School}

\section{Sợ Vì Không Theo Kịp}
\begin{truyen}{Sợ Vì Không Theo Kịp}{Afraid of Falling Behind}
\chuhoa{C}{ó những học sinh không ghét trường. Các em chỉ sợ vào lớp rồi lại thấy mình chậm nhất.}
\textit{(Some students do not hate school. They simply fear entering class and again feeling like the slowest one there.)}

Nỗi sợ này không ồn ào.
\textit{(This fear is not loud.)}

Có em thú thật: “Ngày nào tới cổng trường em cũng phải tự nhủ đừng run.”
\textit{(Some students confess, “Every day at the school gate I have to tell myself not to shake.”)}

Nó nằm trong cái nhìn né tránh, trong việc không dám giơ tay, trong cách ngồi co lại mỗi lần giáo viên gọi tới.
\textit{(It lives in avoiding eye contact, not daring to raise a hand, and shrinking in the seat when the teacher calls.)}

Một nơi học mà luôn làm trẻ thấy mình thua kém sẽ rất khó thành nơi các em muốn tới.
\textit{(A place of learning that constantly makes children feel inferior will struggle to become a place they want to return to.)}

\end{truyen}

\begin{baihoc}
Muốn học sinh bớt sợ đi học, lớp học phải bớt biến sự yếu hơn thành sự xấu hổ.
\textit{(If we want students to fear school less, the classroom must stop turning weakness into shame.)}
\end{baihoc}

\begin{ghinhoanh}
\textbf{Short English to remember:}

Students fear school less when weakness is not turned into shame.

\end{ghinhoanh}

\begin{tuhocnhanh}
\textbf{Gợi ý để dạy và tự nghĩ / Teaching and reflection prompts}

1. Nỗi sợ này biểu hiện ra sao? \textit{How does this fear show itself?}

2. Điều gì trong lớp làm nó lớn lên? \textit{What in the classroom makes it grow?}

3. Mình có đang vô tình biến yếu thành nhục không? \textit{Am I turning weakness into humiliation without noticing?}
\end{tuhocnhanh}
\ngancach

\section{Sợ Vì Bạn Bè Chứ Không Phải Vì Bài Vở}
\begin{truyen}{Sợ Vì Bạn Bè Chứ Không Phải Vì Bài Vở}{Afraid of Classmates More Than Lessons}
\chuhoa{C}{ó học sinh lo buổi sáng không phải vì kiểm tra, mà vì phải gặp lại vài người bạn làm mình co lại.}
\textit{(Some students dread mornings not because of tests, but because they must meet classmates who make them shrink.)}

Trường học với các em không đáng sợ ở bảng, mà ở giờ ra chơi, ở hành lang, ở chỗ ngồi.
\textit{(School is frightening to them not at the board, but in break time, hallways, and seating arrangements.)}

Nếu người lớn chỉ nhìn việc học, ta sẽ bỏ sót nguyên nhân thật của nỗi sợ đi học.
\textit{(If adults look only at academic work, they miss the true source of school fear.)}

Một đứa trẻ không an toàn giữa bạn bè rất khó an tâm để học.
\textit{(A child who feels unsafe among peers struggles to feel calm enough to learn.)}

\end{truyen}

\begin{baihoc}
Muốn cứu việc học của một học sinh, có khi trước hết phải cứu chỗ em đứng giữa bạn bè.
\textit{(Sometimes to save a student's learning, we must first save their place among peers.)}
\end{baihoc}

\begin{ghinhoanh}
\textbf{Short English to remember:}

Sometimes learning improves only after safety among peers is restored.

\end{ghinhoanh}

\begin{tuhocnhanh}
\textbf{Gợi ý để dạy và tự nghĩ / Teaching and reflection prompts}

1. Vì sao nỗi sợ bạn bè làm việc học đi xuống? \textit{Why does peer fear damage learning?}

2. Người dạy cần quan sát thêm ở đâu ngoài tiết học? \textit{Where should teachers look beyond lesson time?}

3. Mình có đang bỏ sót nguồn gốc thật của nỗi sợ này không? \textit{Am I missing the real source of this fear?}
\end{tuhocnhanh}
\ngancach

\section{Sợ Bị Gọi Tên Trước Đám Đông}
\begin{truyen}{Sợ Bị Gọi Tên Trước Đám Đông}{Afraid of Being Called Out Before Everyone}
\chuhoa{C}{ó học sinh vừa nghe tên mình được gọi lớn trước lớp là tim đã đi nhanh.}
\textit{(Some students feel their heart race the moment their name is called out in front of the class.)}

Không phải vì các em ghét chú ý theo mọi nghĩa.
\textit{(It is not necessarily because they hate attention in every sense.)}

Mà vì các em sợ sai trước đám đông, sợ đứng hớ, sợ để lộ chỗ yếu của mình công khai.
\textit{(They fear making mistakes publicly, looking awkward, and exposing weaknesses in front of everyone.)}

Một lớp học nhân hậu cần biết phân biệt giữa khích lệ và đẩy trẻ vào hoảng.
\textit{(A compassionate classroom must know the difference between encouragement and pushing a child into panic.)}

\end{truyen}

\begin{baihoc}
Không phải cơ hội phát biểu nào cũng là cơ hội phát triển nếu nó được đưa bằng cách làm học sinh hoảng sợ.
\textit{(Not every chance to speak is a chance to grow if it is delivered through fear.)}
\end{baihoc}

\begin{ghinhoanh}
\textbf{Short English to remember:}

Encouragement stops being helpful when it pushes a student into panic.

\end{ghinhoanh}

\begin{tuhocnhanh}
\textbf{Gợi ý để dạy và tự nghĩ / Teaching and reflection prompts}

1. Vì sao một số học sinh sợ bị gọi tên? \textit{Why are some students afraid of being called on?}

2. Khích lệ khác ép ở đâu? \textit{How is encouragement different from pressure?}

3. Mình có thể tạo cảm giác an toàn trước khi yêu cầu các em bước ra không? \textit{Can I create safety before asking them to step forward?}
\end{tuhocnhanh}
\ngancach

\section{Em Có Thể Đi Học Mà Không Phải Co Người Lại Không}
\begin{truyen}{Em Có Thể Đi Học Mà Không Phải Co Người Lại Không}{Can I Go to School Without Shrinking Myself}
\chuhoa{C}{ó một ước muốn rất đơn giản của nhiều học sinh: được đi học mà không phải luôn đề phòng.}
\textit{(Many students carry a simple wish: to go to school without constantly bracing themselves.)}

Không phải đề phòng thầy cô, bạn bè, điểm số, hay ánh mắt chê.
\textit{(Not bracing against teachers, classmates, grades, or mocking looks.)}

Khi trường học làm được điều đó, nó mới thật sự là nơi nuôi người.
\textit{(When school can do that, it truly becomes a place that nourishes people.)}

Học tốt là quan trọng, nhưng học trong tư thế co rút mãi thì rất tốn đời sống bên trong.
\textit{(Learning well matters, but learning while always contracted consumes inner life.)}

\end{truyen}

\begin{baihoc}
Trường học tốt không chỉ dạy kiến thức. Nó còn cho học sinh một chỗ để thả lỏng mà vẫn được tôn trọng.
\textit{(A good school does not only teach knowledge. It also gives students a place where they can relax and still be respected.)}
\end{baihoc}

\begin{ghinhoanh}
\textbf{Short English to remember:}

A good school lets students relax without losing respect.

\end{ghinhoanh}

\begin{tuhocnhanh}
\textbf{Gợi ý để dạy và tự nghĩ / Teaching and reflection prompts}

1. Điều gì làm học sinh phải co người lại ở trường? \textit{What makes students shrink themselves at school?}

2. Một nơi học an toàn cần có những dấu hiệu nào? \textit{What signs show that a learning place is safe?}

3. Mình đang giúp lớp học thả lỏng hơn hay co lại hơn? \textit{Am I helping the classroom open up or contract?}
\end{tuhocnhanh}
\ngancach

\section{Em Sợ Chuông Báo Kiểm Tra Hơn Sợ Bài}
\begin{truyen}{Em Sợ Chuông Báo Kiểm Tra Hơn Sợ Bài}{I Fear the Test Announcement More Than the Lesson}
\chuhoa{C}{ó học sinh nghe tới chữ kiểm tra là mặt đổi khác hẳn.}
\textit{(Some students' faces change completely at the word test.)}

Không phải em không học.
\textit{(It is not that they have not studied.)}

Em sợ cảm giác tim đập nhanh, tay lạnh, đầu trống, và ánh mắt người lớn chờ con điểm.
\textit{(They fear the racing heart, cold hands, blank mind, and the adult gaze waiting for a number.)}

Nỗi sợ đi học đôi khi không nằm ở kiến thức, mà nằm ở cơ chế đánh giá bủa quanh kiến thức.
\textit{(Fear of school sometimes lives not in knowledge itself, but in the machinery of evaluation surrounding it.)}

\end{truyen}

\begin{baihoc}
Muốn học sinh bớt sợ học, người lớn cũng phải nhìn lại cách mình làm các em sợ việc bị đo.
\textit{(If adults want students to fear learning less, they must also reconsider how fear is built around being measured.)}
\end{baihoc}

\begin{ghinhoanh}
\textbf{Short English to remember:}

Sometimes students fear the measurement more than the learning.

\end{ghinhoanh}

\begin{tuhocnhanh}
\textbf{Gợi ý để dạy và tự nghĩ / Teaching and reflection prompts}

1. Vì sao chữ “kiểm tra” lại làm học sinh đổi mặt? \textit{Why does the word “test” change a student's face?}

2. Nỗi sợ này nằm ở đâu ngoài kiến thức? \textit{Where does this fear live besides the subject matter?}

3. Mình có đang tạo áp lực đánh giá quá dày không? \textit{Am I building too much evaluative pressure?}
\end{tuhocnhanh}
\ngancach

\section{Em Sợ Đi Học Vì Biết Thế Nào Cũng Bị Gọi Lên}
\begin{truyen}{Em Sợ Đi Học Vì Biết Thế Nào Cũng Bị Gọi Lên}{I Fear School Because I Know I Will Be Called Out}
\chuhoa{C}{ó học sinh vào lớp là ngồi tính trước xem hôm nay mình có bị gọi lên bảng không.}
\textit{(Some students enter class already calculating whether they will be called to the board today.)}

Các em không còn ở trong tư thế học.
\textit{(They are no longer in a position to learn.)}

Các em ở trong tư thế canh chừng.
\textit{(They are in a position of constant vigilance.)}

Một đứa trẻ học trong canh chừng quá lâu sẽ coi trường như nơi phải thủ chứ không phải nơi được lớn.
\textit{(A child who learns under vigilance for too long sees school as a place to defend against, not a place to grow in.)}

\end{truyen}

\begin{baihoc}
Học sinh khó học tốt khi hệ thần kinh của các em đang bận chuẩn bị cho nỗi xấu hổ tiếp theo.
\textit{(Students cannot learn well when their nervous system is busy preparing for the next humiliation.)}
\end{baihoc}

\begin{ghinhoanh}
\textbf{Short English to remember:}

Students struggle to learn when they are busy preparing for shame.

\end{ghinhoanh}

\begin{tuhocnhanh}
\textbf{Gợi ý để dạy và tự nghĩ / Teaching and reflection prompts}

1. Vì sao sự canh chừng làm việc học tê lại? \textit{Why does constant vigilance numb learning?}

2. Học sinh đang chuẩn bị cho điều gì thay vì chuẩn bị học? \textit{What are students preparing for instead of preparing to learn?}

3. Mình có đang để lớp học nghiêng về canh chừng không? \textit{Am I making the classroom feel like a place of vigilance?}
\end{tuhocnhanh}
\ngancach

\section{Có Hôm Em Không Muốn Đến Trường Chỉ Vì Một Cái Bàn}
\begin{truyen}{Có Hôm Em Không Muốn Đến Trường Chỉ Vì Một Cái Bàn}{Some Days One Desk Is Enough to Make Me Fear School}
\chuhoa{C}{ó học sinh sợ đi học chỉ vì biết hôm nay vẫn phải ngồi ở chỗ đó, cạnh người đó, trước ánh mắt đó.}
\textit{(Some students fear school simply because they know they will sit in that seat, next to that person, under that look.)}

Người lớn nghĩ nỗi sợ đi học phải là chuyện lớn.
\textit{(Adults assume school fear must come from something large.)}

Nhưng nhiều khi nó nằm ngay ở một chỗ ngồi rất cụ thể.
\textit{(But often it lives in something as concrete as one desk.)}

Một chi tiết nhỏ lặp lại mỗi ngày có thể biến thành một áp lực rất to.
\textit{(A small detail repeated daily can grow into a very large pressure.)}

\end{truyen}

\begin{baihoc}
Muốn hiểu nỗi sợ đi học, đôi khi phải hỏi học sinh sợ nhất đoạn nào trong một ngày ở trường.
\textit{(To understand fear of school, adults sometimes need to ask which exact part of the school day feels worst.)}
\end{baihoc}

\begin{ghinhoanh}
\textbf{Short English to remember:}

Sometimes one repeated small detail becomes a major source of fear.

\end{ghinhoanh}

\begin{tuhocnhanh}
\textbf{Gợi ý để dạy và tự nghĩ / Teaching and reflection prompts}

1. Vì sao một chi tiết nhỏ lại thành nỗi sợ lớn? \textit{Why can one small detail become a big fear?}

2. Mình có hỏi học sinh đủ cụ thể về điều các em sợ chưa? \textit{Have I asked concretely enough about what students fear?}

3. Lớp học có chi tiết nào đang lặp lại mà làm ai đó khổ? \textit{Is there any repeated detail in the classroom that may be hurting someone?}
\end{tuhocnhanh}
\ngancach

\section{Em Đi Học Với Cái Bụng Căng Như Đi Vào Một Chỗ Bị Xét}
\begin{truyen}{Em Đi Học Với Cái Bụng Căng Như Đi Vào Một Chỗ Bị Xét}{I Go to School Feeling Like I Am Entering a Place of Judgment}
\chuhoa{C}{ó học sinh chuẩn bị đến trường giống như chuẩn bị vào nơi bị soi: tóc tai ổn chưa, đồng phục ổn chưa, bài đủ chưa, hôm nay có bị nhắc gì không.}
\textit{(Some students prepare for school like entering a place of inspection: hair okay, uniform okay, homework okay, will I be criticized today.)}

Đi học mà bụng luôn căng thì rất khó yêu việc học.
\textit{(When the stomach is always tight, it is hard to love learning.)}

Trường học cần kỷ luật.
\textit{(School needs discipline.)}

Nhưng nếu không khí chung nghiêng quá nhiều sang xét nét, học sinh sẽ lớn lên trong tâm thế phòng thủ.
\textit{(But if the atmosphere leans too heavily toward scrutiny, students grow up in a defensive posture.)}

\end{truyen}

\begin{baihoc}
Kỷ luật cần có, nhưng cảm giác chung của trường không nên là cảm giác một tòa án thu nhỏ.
\textit{(Discipline matters, but the general feel of school should not resemble a small courtroom.)}
\end{baihoc}

\begin{ghinhoanh}
\textbf{Short English to remember:}

School needs discipline without feeling like a courtroom.

\end{ghinhoanh}

\begin{tuhocnhanh}
\textbf{Gợi ý để dạy và tự nghĩ / Teaching and reflection prompts}

1. Vì sao cảm giác bị xét làm trẻ mệt? \textit{Why does feeling judged exhaust children?}

2. Kỷ luật khác xét nét ở đâu? \textit{How is discipline different from constant scrutiny?}

3. Trường mình đang nghiêng về bên nào hơn? \textit{Which side is our school leaning toward?}
\end{tuhocnhanh}
\ngancach

\section{Một Câu “Mai Em Vẫn Cứ Tới Nhé” Có Khi Đỡ Được Nỗi Sợ Đi Học}
\begin{truyen}{Một Câu “Mai Em Vẫn Cứ Tới Nhé” Có Khi Đỡ Được Nỗi Sợ Đi Học}{One Sentence Can Lessen Fear of School}
\chuhoa{C}{ó hôm một em làm sai, mặt tái đi vì nghĩ mai chắc không dám vào lớp.}
\textit{(One day a student made a mistake and turned pale, convinced they would not dare come back tomorrow.)}

Tôi chỉ nói: “Mai em vẫn cứ tới nhé, chuyện này mình xử tiếp.”
\textit{(I only said, “Come tomorrow as usual, we'll continue dealing with this then.”)}

Em đứng yên một lúc rồi gật.
\textit{(The student stood still for a moment and nodded.)}

Có những nỗi sợ đi học được giữ xuống chỉ nhờ biết mai mình vẫn còn chỗ.
\textit{(Some school fear is held down simply by knowing there is still a place for you tomorrow.)}

\end{truyen}

\begin{baihoc}
Học sinh bớt sợ trường hơn khi các em biết lỗi lầm không làm mình mất luôn quyền quay lại ngày mai.
\textit{(Students fear school less when they know mistakes do not erase their right to return tomorrow.)}
\end{baihoc}

\begin{ghinhoanh}
\textbf{Short English to remember:}

Students fear school less when they know they still have a place tomorrow.

\end{ghinhoanh}

\begin{tuhocnhanh}
\textbf{Gợi ý để dạy và tự nghĩ / Teaching and reflection prompts}

1. Vì sao quyền được quay lại ngày mai quan trọng với học sinh? \textit{Why does the promise of tomorrow matter so much?}

2. Một câu trấn an đúng lúc có thể làm gì? \textit{What can one timely reassuring sentence do?}

3. Mình có biết giữ học sinh lại với trường bằng những câu như vậy không? \textit{Do I know how to keep students connected to school through words like this?}
\end{tuhocnhanh}
\ngancach

\section{Em Muốn Đi Học Mà Không Phải Dũng Cảm Quá Nhiều Mỗi Ngày}
\begin{truyen}{Em Muốn Đi Học Mà Không Phải Dũng Cảm Quá Nhiều Mỗi Ngày}{I Want School Without Needing So Much Courage Every Day}
\chuhoa{C}{ó những học sinh ngày nào đi học cũng phải lấy đà rất mạnh.}
\textit{(Some students need a great deal of courage every single day just to come to school.)}

Người ngoài không thấy vì em vẫn tới đều.
\textit{(Outsiders do not see it because the student still shows up.)}

Nhưng sống kiểu ngày nào cũng phải tự kéo mình qua nỗi sợ là rất hao.
\textit{(But living in a mode where you drag yourself through fear every day is exhausting.)}

Một trường học tốt không bắt học sinh phải dũng cảm quá mức chỉ để hiện diện.
\textit{(A good school does not require excessive courage merely for a student to be present.)}

\end{truyen}

\begin{baihoc}
Nếu việc tới trường ngày nào cũng là một màn gồng lớn, người lớn cần xem lại môi trường chứ không chỉ xem lại học sinh.
\textit{(If coming to school every day feels like a huge act of strain, adults need to review the environment, not only the student.)}
\end{baihoc}

\begin{ghinhoanh}
\textbf{Short English to remember:}

A good school should not require extreme courage just to attend.

\end{ghinhoanh}

\begin{tuhocnhanh}
\textbf{Gợi ý để dạy và tự nghĩ / Teaching and reflection prompts}

1. Vì sao đi học đều chưa chắc là đi học nhẹ nhàng? \textit{Why does regular attendance not mean easy attendance?}

2. Môi trường có thể đang buộc học sinh phải gồng ở đâu? \textit{Where might the environment be forcing students to strain?}

3. Mình có đang xem lại lớp học đủ nghiêm túc như xem lại học sinh không? \textit{Am I examining the classroom as seriously as I examine the student?}
\end{tuhocnhanh}
